\documentclass[11pt]{article}
\usepackage[margin=1in]{geometry}
\usepackage{amsmath,amssymb,amsthm}
\usepackage[colorlinks=true,urlcolor=blue]{hyperref}
\newtheorem{theorem}{Theorem}
\newtheorem{lemma}{Lemma}
\newtheorem{corollary}{Corollary}
\newcommand{\R}{\mathbb R}
\newcommand{\pdet}{\operatorname{pdet}}
\newcommand{\vol}{\operatorname{vol}}
\newcommand{\Tr}{\operatorname{tr}}
\title{A degree-aware one-edge budget for fractional Markov forms}
\author{Working research note; classical ingredients, priority unverified}
\date{November 15, 2026}
\begin{document}
\maketitle

\paragraph{Scope.}
This note supplies a rigorously derived extension of the repository's
one-edge fractional-response budget to an energy form with changing vertex
degrees. It is not a new Cheeger inequality, a partitioning algorithm, a
sharp extremal count theorem, or a claim of historical novelty. Joint
monotonicity is a classical operator-mean property; the proof below is
included to settle the singular/common-kernel issue directly.

Let $G$ be a finite connected loopless undirected graph with nonnegative
edge weights, $n\ge2$, positive degrees $d_v$, degree matrix $D$, and
Laplacian $L$. Its normalized Laplacian is
$M=D^{-1/2}LD^{-1/2}$. Fix $0<\alpha<1$ and define the symmetric
fractional Markov energy form
\[
 H_\alpha(G)=D^{1/2}M^\alpha D^{1/2}.
\]
Equivalently, $H_\alpha=D\#_\alpha L$, the weighted operator geometric
mean, with its continuous extension to positive semidefinite second input.
It is the energy form of $(D^{-1}L)^\alpha$ in $L^2(D)$.
Its kernel is the fixed space $\R\mathbf1$, although the kernel of $M$
changes when the degrees change.

Insert an edge of weight $w>0$ between distinct $u,v$, or increase an
existing edge weight by $w$. Put
\[
 b=e_u-e_v,\quad L'=L+wbb^T,\quad
 D'=D+w(e_ue_u^T+e_ve_v^T),\quad V=\sum_xd_x,
 \quad \tau=w b^TL^\dagger b.
\]
Write $H=H_\alpha(G)$ and $H'=H_\alpha(G')$.

\begin{theorem}[Exact degree-aware logdet budget]
\label{thm:budget}
The response is positive semidefinite: $H'\succeq H$.
On $\mathbf1^\perp$, let
\[
 C=H^{-1/2}(H'-H)H^{-1/2},
\]
where this expression uses the inverse of the restriction of $H$ to
$\mathbf1^\perp$ and is extended by zero on constants. Then $C\succeq0$
and
\begin{align*}
 \log\det_{\mathbf1^\perp}(I+C)
 &=B_\alpha(G;uv,w)\\
 &:=\alpha\log(1+\tau)
 +(1-\alpha)\log\!\left(
 \frac{(1+w/d_u)(1+w/d_v)}{1+2w/V}\right).
\end{align*}
In particular,
\[
 0\le B_\alpha(G;uv,w)\le(2-\alpha)\log(1+\tau).
\]
For $\varepsilon>0$,
\[
 \#\{j:\lambda_j(C)>\varepsilon\}
 <\frac{B_\alpha(G;uv,w)}{\log(1+\varepsilon)}.
\]
The last display is an upper bound only; sharpness is not asserted.
\end{theorem}

\begin{proof}[Proof of positivity]
For $A\succ0$, $B\succeq0$, their parallel sum satisfies
\[
 x^T(A:B)x=\inf_{y+z=x}(y^TAy+z^TBz)
 =x^T[A-A(A+B)^{-1}A]x.
\]
The variational expression proves joint monotonicity in $A,B$, including
singular $B$. The scalar fractional-power integral, applied to $M$, gives
\[
 H_\alpha(G)
 =\frac{\sin(\pi\alpha)}\pi
 \int_0^\infty t^{\alpha-1}
 \bigl[D-tD(L+tD)^{-1}D\bigr],dt
 =\frac{\sin(\pi\alpha)}\pi
 \int_0^\infty t^{\alpha-1}\bigl[D:(L/t)\bigr],dt.
\]
The integral converges: the bracket is bounded by $D$ near zero and by
$L/t$ near infinity. Both $D'\succeq D$ and $L'\succeq L$, so joint
monotonicity gives the desired order at every $t$, and hence after
integration. The common kernel is immediate from $\ker M=\R D^{1/2}\mathbf1$.
\end{proof}

\begin{lemma}[Pseudodeterminant under congruence]
\label{lem:pdet}
If $A\succeq0$ has rank $n-1$ and unit kernel vector $z$, and $S$ is
invertible, then
\[
 \pdet(SAS^T)=\det(S)^2\pdet(A)\,\|S^{-T}z\|^2.
\]
\end{lemma}
\begin{proof}
Compare the coefficients of $t$ in
\[
 \det(SAS^T+tI)=\det(S)^2\det(A+tS^{-1}S^{-T}).
\]
For a rank-$n-1$ positive matrix,
$\operatorname{adj}(A)=\pdet(A)zz^T$; thus the coefficient on the right
is $\det(S)^2\pdet(A)z^TS^{-1}S^{-T}z$. The coefficient on the left is
$\pdet(SAS^T)$.
\end{proof}

\begin{proof}[Proof of the exact budget and the count]
Let $T(G)$ denote the sum of products of edge weights over spanning trees.
The weighted matrix-tree theorem gives $\pdet(L)=nT(G)$.
Apply Lemma~\ref{lem:pdet} first to $M=D^{-1/2}LD^{-1/2}$, using the
unit kernel vector $\mathbf1/\sqrt n$ of $L$, to obtain
\[
 \pdet(M)=\frac{VT(G)}{\prod_x d_x}.
\]
Apply it a second time to $H=D^{1/2}M^\alpha D^{1/2}$, using the unit
kernel vector $D^{1/2}\mathbf1/\sqrt V$ of $M^\alpha$. This gives
\[
 \boxed{\displaystyle
 \pdet(H_\alpha(G))=
 n\,T(G)^\alpha\left(\frac{\prod_xd_x}{V}\right)^{1-\alpha}.}
\]
On the fixed support $\mathbf1^\perp$, the determinant lemma yields
$T(G')/T(G)=1+wb^TL^\dagger b=1+\tau$. The two modified degree factors
and $V'=V+2w$ give precisely the formula for $B_\alpha$ in the theorem.
Because $C\succeq0$, $B_\alpha\ge0$.

For the clean upper bound, effective conductance from $u$ to $v$ is at
most $d_u$: the voltage test equal to one at $u$ and zero at every other
vertex has energy $d_u$. Hence $R_{uv}=b^TL^\dagger b\ge1/d_u$;
likewise $R_{uv}\ge1/d_v$. Consequently $w/d_u,w/d_v\le\tau$, and
the denominator $1+2w/V\ge1$ can be discarded for an upper bound.
This gives $B_\alpha\le[\alpha+2(1-\alpha)]\log(1+\tau)$.
Finally, if $m$ eigenvalues of $C$ strictly exceed $\varepsilon$, then
\[
 m\log(1+\varepsilon)
 <\sum_j\log(1+\lambda_j(C))=B_\alpha.
\]
\end{proof}

\paragraph{Two checks.}
For a two-vertex graph with edge weight $c$, the only positive eigenvalue
of $H_\alpha$ is $c2^\alpha$. Inserting additional weight $w$ gives
$B_\alpha=\log(1+w/c)$ for every $\alpha$, agreeing with the theorem.
As $\alpha\uparrow1$, the formula reduces to the usual rank-one
Laplacian determinant budget $\log(1+\tau)$. Neither check proves
sharpness of the count for fixed $0<\alpha<1$.

\paragraph{Why this does not yet solve partitioning.}
For every $f$, writing $x=D^{1/2}f$ gives
\[
 \frac{f^TH_\alpha f}{f^TDf}
 =\frac{x^TM^\alpha x}{x^Tx}.
\]
On $\ker M^\perp$, concavity gives an upper bound by
$[x^TMx/(x^Tx)]^\alpha$. It does not give a uniform reverse comparison
between fractional energy and original energy near spectral zero.
For example an eigenvector of $M$ at eigenvalue $\lambda>0$ has ratio
$\lambda^\alpha/\lambda=\lambda^{\alpha-1}$, which diverges as
$\lambda\downarrow0$. Thus a gap-free relative approximation to
$H_\alpha$ is not automatically a relative approximation to $L$, nor a
rounding or balanced-cut guarantee. The separately changing mass matrix
$D$ must also be retained in any generalized spectral calculation.

\begin{corollary}[No gap-free conversion of the permitted residual]
Fix $0<\alpha<1$ and $\varepsilon>0$. There is no finite constant
$K_{\alpha,\varepsilon}$ such that, for every connected simple unweighted
graph and every positive semidefinite residual $E$ satisfying
$0\preceq E\preceq\varepsilon H_\alpha(G)$, one necessarily has
$E\preceq K_{\alpha,\varepsilon}L$.
This failure already occurs for degree-two regular graphs and rank-one $E$.
\end{corollary}
\begin{proof}
Take the cycle $C_n$, so $D=2I$ and
$\lambda_n=1-\cos(2\pi/n)>0$ is the first positive eigenvalue of $M$.
Choose a unit eigenvector $v_n$ at $\lambda_n$ and put
$E_n=2\varepsilon\lambda_n^\alpha v_nv_n^T$.
Then $0\preceq E_n\preceq\varepsilon H_\alpha(C_n)$, but the
smallest $K$ such that $E_n\preceq K L(C_n)$ is
$\varepsilon\lambda_n^{\alpha-1}\to\infty$.
\end{proof}
The corollary concerns the generic residual guarantee alone. It does not
preclude sharper graph-specific rounding analysis or a smaller structured
residual produced by a particular algorithm.

\paragraph{Classical provenance.}
The parallel sum and its order properties originate in
W.~N.~Anderson, Jr. and R.~J.~Duffin,
\emph{Series and parallel addition of matrices},
J.~Math.~Anal.~Appl.~26 (1969), 576--594,
\url{https://doi.org/10.1016/0022-247X(69)90200-5}.
The operator-mean framework is F.~Kubo and T.~Ando,
\emph{Means of positive linear operators},
Math.~Ann.~246 (1980), 205--224,
\url{https://doi.org/10.1007/BF01371042}.
The determinant computation is a direct combination of these ideas with
the classical weighted matrix-tree theorem and determinant lemma.
No claim is made that this particular graph formula is historically new.
\end{document}
